% Proposal from round-1 reviewer opus_contribution. NOT integrated; the coordinator owns manuscript edits.
% Architecture A: label validity. Replaces the abstract's thesis, Introduction paragraphs 3-6, and the
% G13 paragraph in Section "Social conduct". Numbers marked [R] come from this reviewer's read-only
% script (scripts/end_provenance.py) and need coordinator re-verification before use. Other numbers are
% the draft's retained figures. Uses existing preamble macros (\decision, \game).

% ---------------------------------------------------------------- Title
% What Agent Logs Do Not Say: Resolution, Reachability and Delivery in
% Long-Horizon Multi-Agent Economic Competition

% ---------------------------------------------------------------- Abstract (replacement)
Evaluations of long-horizon agents routinely read behavioral labels off logs: an agent
declined to invest, could have paid its debt, made an offer, informed a rival, or was
persuaded. We show that such labels fail in a structured way, and we identify where. In
twenty complete four-player games of MonopolyBench (6,989 decisions, 8,679 response
attempts), three layers separate model text from state change. The response-resolution
layer decides whether an applied action is model-authored. The accessibility layer
separates displayed menus and nominal wealth from reachable legal sequences. The
delivery layer separates canonical speech and harness summaries from what a recipient
received. Of 663 applied end-of-turn actions on menus where construction was legal, 90
were not model-chosen holds; for one configuration the figure is 61 of 104, against
under 2\% for two others. All 61 of that configuration's defaulted construction plans
execute when the same purchases are reordered. Of 60 eliminations, 56 are insolvent
under a generous nominal bound, but every one of the four remaining cases changes label
under legal-sequence execution. A harness summary reversed giving directions in 44 of
113 accepted trades and reached 2,054 prompts. In social exchanges, settlements occur
while recipients correct the false premise, and a received offer can contradict the
proposer's own accompanying note. Exact-prefix engine execution adjudicates feasibility
and exposure; it does not identify persuasion or intent, which require re-sampling the
recipient policy at a fixed prefix. Behavioral rates computed without these layers are
partly measurements of the harness.

% ---------------------------------------------------------------- Introduction core (replaces paras 3-6)
\paragraph{The problem.}
A growing literature measures how language agents invest, bargain, deceive and keep
commitments over long interactions, and reports rates: of misaligned messages, of promises
enacted, of trades accepted, of avoidable failures~\cite{li2026arena,wu2026agentmonopoly,hu2026finboardbench}.
Each rate presupposes a label that maps a logged record to a behavior. An END action is
read as a choice to hold, a bankruptcy declaration as surrender of a solvent position, a
sent offer as the proposer's intent, a logged message as information the recipient
had, and a closed deal as evidence of persuasion. These mappings are rarely tested,
because testing them requires knowing which response produced the applied action, which
legal sequences were reachable from the exact state, and what each recipient's prompt
actually contained.

\paragraph{Three layers.}
MonopolyBench records all three. The \emph{resolution} layer comprises refusals,
corrective retries, recovery prompts, deterministic fallbacks and forced tool choice; it
determines whether an applied action is model-authored and under which prompt condition.
The \emph{accessibility} layer is the gap between what a menu displays or a balance sheet
implies and what a legal sequence of the agent's own actions can reach, under rules such
as even selling and finite building stock. The \emph{delivery} layer is the gap between
canonical events and each recipient's bounded memory and harness-written summaries. We
restore exact historical prefixes and execute specified legal sequences through the
engine, and we reconstruct every delivered memory window, to adjudicate labels against
these layers.

\paragraph{What we find.}
Naive labels are usually right on average and often wrong where analysts look. Of the
663 END actions applied while construction was legal, 573 are model-chosen by provenance;
the remaining 90 are engine defaults or received tool calls whose action contradicts the
accompanying content [R]. The error is concentrated by configuration: 61 of 104 for a
participant whose provider refused its standard prompt in 1,605 decisions, against four
of 236 and one of 144 for two others [R]. A cross-model ``propensity to hold cash'' would
therefore rank harness conditions. Among sixty eliminations, the nominal-value label is
correct for the 56 that no own-asset liquidation could rescue. It changes for all four
remaining cases: three admit a legal current-debt payment, while one has \$1,310 in nominal
value but only \$440 accessible because a hotel cannot be downgraded. In the delivery
layer, one summary assumption reversed giving directions in 44 accepted trades and
reached 2,054 prompts, and 149 nonempty public messages never entered any timeline.
Finally, our own analysis produced fluent, mutually agreed narratives that the record
contradicted. These include a misdated productive group, a single construction step
mistaken for cumulative spending, and a menu field read as a prohibition.

\paragraph{Social claims as a stress test.}
Social exchanges show the limits of the instruments as clearly as their reach. In two
settled trades the recipient states the correct proposition while accepting, so
settlement does not imply uptake of a false premise. In another, the received proposal
carries a private note instructing that it not be sent. The raw provider response,
which was not streamed, contains both the proposal and a contradictory end-of-turn call.
So ``sent'' does not imply ``intended'', and the divergence originates upstream of our
harness. These are counterexamples to sufficiency rules, and one checked counterexample
suffices to reject a rule. They are not prevalence estimates. Where the question is
whether a message \emph{changed} a recipient's choice, engine branches are the wrong
instrument: they can show that a seller's implied constraint was false, not that a buyer
would have paid less without hearing it. That question requires re-querying the
recipient policy at the identical prefix with the message ablated, which the validated
prefix-restoration machinery supports and which we specify as the next experiment.

\paragraph{Contributions.}
(i) A decomposition of log-derived behavioral labels into resolution, accessibility and
delivery layers, with full-denominator naive-versus-adjudicated comparisons in a cohort
of complete long-horizon games. (ii) An exact-prefix legal-sequence procedure that
adjudicates feasibility claims made by agents and by analysts, together with an explicit
account of the claims it cannot adjudicate. (iii) Evidence that label error is
configuration-specific. Cross-model behavioral comparisons should therefore report
action provenance and information condition before reporting rates. The cohort supports
no model ranking, capability trend or misalignment rate. Public deception was permitted,
and capability was not varied. We do not claim that any other benchmark exhibits our
specific error rates. We claim only that the layers exist in any harness with retries,
fallbacks, tool calling and bounded memory, and that their contribution is measurable
when the record is complete.

% ---------------------------------------------------------------- G13 paragraph (replacement)
\paragraph{\game{13}: an omission whose premise both parties shared, and whose purpose failed.}
From \decision{13}{218} Grok (seat~4) repeatedly offers mortgaged Pennsylvania Avenue,
first to Sol, arguing that its bankruptcy would transfer the yellows. Sol publicly
corrects one such claim at \decision{13}{244}, and both messages are in Opus's delivered
timeline at \decision{13}{251}. At \decision{13}{250} Grok makes the same appeal to Opus,
misstating the Short Line risk as 5/36; from Atlantic it is 4/36. Opus counters at \$130,
and its returned reasoning notes that the payment ``keeps p4 solvent enough to cover Short
Line rent to p2''. At \decision{13}{252} Grok says that \$130 ``leaves me \$35 short''. The
cash figure is true. A Ventnor mortgage would reach \$295, as an exact-prefix branch
confirms. Grok's applied note acknowledges that route, and its returned reasoning plans to
omit it as ``slightly deceptive'', betting that the buyer ``may not calculate the
mortgage''. The omitted fact is the seller's own outside option, not a false assertion
about the board. At \decision{13}{253} Opus accepts \$200, reasoning that ``player~2 would
inherit the yellows plus Pennsylvania''. This nearly reproduces the d252 message, so
uptake of the seller's consequence framing is evidenced. Uptake of the false minimum is
neither evidenced nor excluded. The shared premise did not hold for long. At
\decision{13}{311}, 43 turns later, Grok lands on Short Line, declares bankruptcy, and Sol
inherits the yellows. As with the denial purchase in game~18, buying a rival's solvency
under creditor inheritance lasts only as long as that solvency.
